
\section{Results}

\subsection{Data}
The 90 FAERS extracts held 24{,}812{,}424 report versions, from which 20{,}692{,}682
cases remained after keeping the latest version; 104{,}186 cases on FDA deletion lists, cases without a
reaction or mappable ingredient, 1{,}681 report identifiers recorded under two cases and
510{,}312 content duplicates were then removed, leaving \textbf{19{,}876{,}330
reports} (4\% of drug rows could not be mapped; Supplementary Figure~S1). They contained
83{,}883 hearing-loss and 42{,}900 tinnitus reports. Hearing-loss
reports came more often from consumers and from patients aged $\geq$65 years than other reports (Supplementary Table~S4).
The primary adjusted population excluded 19{,}277 reports with an ear or hearing
indication, 1{,}842{,}774 reports naming a listed ototoxic drug and
486{,}359 lawyer reports, leaving 17{,}571{,}385 reports
(70{,}438 hearing-loss cases). Canada Vigilance contributed 1{,}210{,}857 reports
after removing 58{,}623 duplicates (4{,}000 hearing-loss and
3{,}284 tinnitus reports).

Drugs of the established ototoxic classes gave strong signals in both databases (Supplementary Table~S2), for example
amikacin (FAERS adjusted OR 12.7 (11.6--13.9)), gentamicin (11.4 (9.8--13.2)) and cisplatin (4.3 (4.0--4.6)); so did
the PDE5 inhibitors sildenafil (3.9 (3.7--4.2)) and tadalafil (4.2 (3.9--4.6)), whose labels carry hearing-loss
warnings, and teprotumumab (59.8 (55.4--64.5)). Several label-positive drugs did not signal: oxaliplatin (crude ROR
0.44), vardenafil (6 cases), and, after adjustment,
carboplatin (crude 1.96, adjusted 0.82) and
vincristine, whose crude associations were carried largely by reports that also named another listed drug---the
sensitivity that the co-medication filter costs.

\subsection{Adjustment changes the signal list but not its agreement with labels}
In FAERS without Canadian reports the crude rule flagged 300 ingredients for hearing loss and the adjusted
rule 310; 63 crude signals were lost and 73 appeared
(Table~\ref{tab:methods}). Against the reference set (106 label-positive and 927 label-negative
ingredients with $\geq$1{,}000 suspect reports), discrimination was moderate and practically identical for the crude and
adjusted lower bounds (AUC 0.718 and 0.723; difference
0.005, 95\% CI -0.023 to 0.030, which excludes an
improvement larger than 0.03). The same held for point estimates
(0.700 vs 0.705), for the IC$_{025}$
(0.720), at lower report thresholds, without the 29 listed drugs, and for tinnitus
(0.775 vs 0.777). Adjustment lowered sensitivity
(0.40 to 0.35) with little change in specificity (0.866 to 0.873).

\subsection{Replication in Canada Vigilance}
Under the pre-specified rule, adjusted hearing-loss signals already on a label replicated in 34/50 (68\%)
of testable cases, unlabelled adjusted signals in 31/76 (41\%), label-negative adjusted signals in
21/54 (39\%), crude-only signals in 9/21 (43\%), and the background of label-negative
non-signals in 43/214 (20\%) (Table~\ref{tab:methods}). The power-calibrated comparison separated the groups
more clearly (Figure~\ref{fig:repl}): labelled signals replicated exactly as often as their FAERS strength predicted
(31 observed vs 31.2 expected (O/E 1.00)), whereas unlabelled signals replicated less than half as often
(25 observed vs 59.5 expected (O/E 0.42)), like label-negative signals (18 observed vs 39.0 expected (O/E 0.46)). Tinnitus showed
the same contrast (60 observed vs 64.4 expected (O/E 0.93) for labelled and 22 observed vs 49.4 expected (O/E 0.45) for unlabelled
signals).

The background replicated more often than chance (36 observed vs
8.5 expected without any association). Solicited reports explained much of this:
\emph{Study} reports, which include manufacturers' patient-support programmes, made up 61\%
of Canadian hearing-loss cases for the background drugs that replicated, against 14\% for
replicated labelled signals and 36\% overall; for ocrelizumab, mepolizumab and omalizumab the share
was 93\%, 93\% and 86\%, against
1\% for gentamicin and 1\% for cisplatin. The mirror image appeared in FAERS:
tafamidis (adjusted OR 4.3 (3.8--4.8)) had
361 hearing-loss reports in FAERS but none in Canada Vigilance against 3.4 expected, and
sacubitril (2{,}801 FAERS reports) none against 1.9 expected (too few to be testable under the
pre-specified rule).

Co-medication bias also replicated when crude estimates were compared. Etoposide, given in cisplatin-containing regimens,
was a crude signal in both databases (FAERS ROR 2.09 (1.91--2.28);
Canada 5.13 (3.68--7.14)) and disappeared after adjustment in both (FAERS
0.67 (0.46--0.98); Canada
1.15 (0.29--4.64)). Temporal replication within FAERS did not separate crude-only
from adjusted hearing-loss signals (early adjusted signals recurred after 2014 in 81/155 (52\%), crude-only
signals in 18/28 (64\%), background 137/938 (15\%)); for tinnitus adjusted signals recurred more
often (100/121 (83\%) vs 9/15 (60\%)).

\subsection{Post-hoc analyses}
With the extended indication filter, Canada Vigilance restricted to unsolicited reports and both estimates adjusted
(Table~\ref{tab:posthoc}), labelled hearing-loss signals replicated in 26/44 (59\%), unlabelled signals in
10/51 (20\%) and the background in 19/158 (12\%) (unlabelled vs background, Fisher
$p=0.24$); calibrated O/E ratios were 0.81
for labelled and 0.20 for unlabelled signals. Handling co-medication by
stratification instead of exclusion, with a short list of strongly ototoxic drugs, or after removing reports listing
$\geq$15 drugs left this contrast unchanged (unlabelled O/E 0.20--0.27
across variants, labelled 0.81--1.02),
as did restricting Canada Vigilance to reports from 2004.

\subsection{Signals from previous screens}
Of the 23 hearing-loss signals that Li et al.\ classified as unlabelled \cite{li2024},
16 were crude and 16 adjusted signals in FAERS without Canadian reports,
6 were testable in Canada and 2 replicated; under the post-hoc
definition 1 remained (glycopyrrolate).
Their strongest unlabelled signals included sacubitril/valsartan and tafamidis. Of the 28 unlabelled
paediatric signals of Liu et al.\ \cite{liu2025}, 11 were adjusted signals in our all-age data and
1 met the post-hoc definition (azithromycin,
which is on the a-priori ototoxic list); cefixime, ceftriaxone, etoposide, isoniazid were lost on adjustment
(Table~\ref{tab:prior}). Only abstract-level lists were available for the remaining screens (Supplementary Table~S6).

\subsection{Case-level review}
Under the post-hoc definition, 10 hearing-loss and 12 tinnitus signals were robust and absent from every US
label (Table~\ref{tab:cand}, Figure~\ref{fig:cand}). The review attributed 4 to report clusters (for
example candesartan--tinnitus: 80\% of cases from the United Kingdom, median 42 drugs per report),
6 to stimulated or mass-exposure reporting (mRNA COVID-19 vaccines, rofecoxib, finasteride),
7 to an indication itself associated with the event (rheumatoid arthritis for etanercept and
tofacitinib; COPD for glycopyrrolate; glaucoma for travoprost; diabetes for insulin aspart; migraine for galcanezumab;
hypothyroidism for liothyronine) and 1 to an established ototoxic class (oral paromomycin).
Four remained without a dominant artefactual explanation: pregabalin and amlodipine for hearing loss, and doxycycline and
chlorthalidone (10 cases) for tinnitus.
The independent assessor agreed on 19 of 22 assignments (Cohen's $\kappa=0.82$);
it placed Tozinameran (hearing loss) in K rather than C; Amlodipine (hearing loss) in D rather than E; Galcanezumab (tinnitus) in E rather than D
(Supplementary Table~S3), which would leave the number of unexplained signals at
4.
